% !TEX program = xelatex
% 2027 新领军学科能力二试 · 语言 · 语文（预测训练卷）
\documentclass[cjk=noto, seal=first]{mojing-exam}

\examtitle{2027 新领军学科能力测试 · 预测训练卷}
\secret{训练用卷 ★ 非官方试题}

\begin{document}

\begin{paper}[time={建议 2 小时（训练设置）}, full=100]{二试 · 语言 · 语文}

\begin{instructions}
  \item 本卷共 4 部分 10 题，满分 100 分，建议用时 2 小时。2027 年简章列出二试“语言（语文、英语）”，未公布题型、时长与分值；本卷全部安排为训练设置。
  \item 现代文阅读材料为本卷编写；文言文材料据通行本整理，有删节，个别字各本有异文。
  \item 书写清楚，作文不少于 800 字。
\end{instructions}

\begin{gradingnotes}
  \item 本评分标准为训练用建议标准。主观题按要点给分，意思相近即可；表述准确、条理清楚者在本档内从高给分。
  \item 作文按内容、表达、发展等级综合评定，见第 10 题说明。
\end{gradingnotes}

\scoretable{一,二,三,四}

% -----------------------------------------------------------------------------
\examsection*{现代文阅读}{本题共 4 小题，20 分}

阅读下面的文章，完成 1—4 题。

\begin{center}\zihao{-4}\heiti 在“懂了”之前\end{center}

\noindent\hspace{2em}很多人都有这样的经验：一道题看了答案，每一步都觉得顺理成章，合上书却写不出来。我们说自己“看懂了”，其实只是“跟上了”。
跟上，是让别人的思路在自己脑中走一遍；懂，是知道这条路为什么非走不可，换一个起点还能自己找到它。

\hspace{2em}数学史上有一个常被讲起的故事。古希腊人发现正方形的对角线与边长不可公度，这个结论几句话就能写完，它带来的震动却延续了很久：
人们不得不重新追问“数”是什么，“长度”又是什么。结论本身并不长，难的是承认它，并且弄清它推翻了什么、保留了什么。一个人若只把它记成
“根号二是无理数”，就等于把一场地震记成了一条天气预报。

\hspace{2em}理解之所以慢，是因为它要求我们在头脑里重建对象之间的关系，而不是把关系的名称存进记忆。名称是可以借来的：我们可以说
“由某某定理可知”，却未必知道这个定理在什么条件下会失效，也未必能在一个陌生的问题里认出它的影子。真正的理解往往伴随着一种“必然感”
——不是因为书上这样写，而是因为换一种写法就会出错。为了得到这种必然感，人们需要亲手试错：把条件去掉一个，看结论还剩下什么；
把例子换一个，看论证在哪一步卡住。

\hspace{2em}有人担心，这样学太慢，赶不上考试的节奏。这个担心并非没有道理。考试有时间限制，熟练确实重要；不熟练的理解像一把没开刃的刀。
但熟练如果不长在理解上，便只是对特定题型的条件反射，题目稍一变形就会失灵。更要紧的是，一个人在学习中养成的，不只是某些知识，
还有面对未知时的态度：是急于找到一个可以套用的模板，还是愿意先停下来，弄清楚自己究竟不懂在哪里。

\hspace{2em}所以，慢一点未必是坏事。在“懂了”之前多停留一会儿，承认困惑，追问原因，把一句结论还原成一连串不得不然的步骤——这段时间看起来没有产出，
却是知识真正扎根的时候。快，可以靠练习获得；而决定一个人能走多远的，往往是他在慢处所下的功夫。

\begin{questions}
\question[3] 下列对文章内容的理解，\textbf{不正确}的一项是 \paren[C]
\choices{作者认为“看懂了”往往只是“跟上了”别人的思路}
        {作者认为记住定理的名称并不等于理解了它的适用条件}
        {作者认为熟练与理解相互排斥，追求熟练会损害理解}
        {作者认为理解中的“必然感”需要通过亲手试错获得}
\begin{solution}
C。第四段说“熟练确实重要”，只是强调熟练要“长在理解上”，并未说二者相互排斥。
\end{solution}

\question[5] 第二段把只记住“根号二是无理数”比作“把一场地震记成了一条天气预报”。这个比喻表达了什么意思？
\answerspace{3cm}
\begin{solution}
结论的表述很短，引起的却是关于“数”与“长度”的根本变革；\score{2}
只记结论，就丢掉了它推翻了什么、保留了什么，把深刻的意义变成了一条轻描淡写的信息；\score{4}
“地震”与“天气预报”的反差突出了这种失落，意在提醒读者理解结论背后的来龙去脉。\score{5}
\end{solution}

\question[6] 根据全文，概括作者所说的“理解”包含哪些要素。
\answerspace{3cm}
\begin{solution}
答出以下任意四点，每点 1.5 分：
① 知道一条思路为什么非走不可（必然性）；② 换一个起点、在陌生问题中仍能自己找到或认出它（迁移）；
③ 清楚结论在什么条件下成立、何时失效（边界）；④ 在头脑中重建对象之间的关系，而不是记住关系的名称；
⑤ 需要通过亲手试错（去掉条件、更换例子）来获得。\score{6}
\end{solution}

\question[6] 第四段是怎样回应“这样学太慢”这一担心的？请分析其论证思路。
\answerspace{3cm}
\begin{solution}
先让步，承认担心有道理：考试有时限，熟练重要，并用“没开刃的刀”作比；\score{2}
再转折，指出不长在理解上的熟练只是条件反射，题目一变形就失灵；\score{4}
最后递进，把讨论从知识提升到“面对未知时的态度”。让步—转折—递进，使论证更周全，也更有说服力。\score{6}
\end{solution}
\end{questions}

% -----------------------------------------------------------------------------
\examsection*{文言文阅读}{本题共 4 小题，20 分}

阅读下面的材料，完成 5—8 题。

\textbf{材料一}　（《九章算术·方田》“半周半径相乘得积步”，刘徽注，有删节）

{\kaishu\hspace{2em}以六觚之一面乘半径，因而三之，得十二觚之幂。若又割之，次以十二觚之一面乘半径，因而六之，则得二十四觚之幂。
割之弥细，所失弥少。割之又割，以至于不可割，则与圆周合体而无所失矣。觚面之外，又有余径。以面乘余径，则幂出弧表。若夫觚之细者，
与圆合体，则表无余径。表无余径，则幂不外出矣。以一面乘半径，觚而裁之，每辄自倍。故以半周乘半径而为圆幂。\par}

\textbf{材料二}　（《墨子·经上》，节选）

{\kaishu\hspace{2em}平，同高也。\quad 中，同长也。\quad 直，参也。\quad 圜，一中同长也。\quad 方，柱隅四讙也。\quad 倍，为二也。\par}

\smallskip
\noindent{\small 注：觚（gū），多边形的角，此处指正多边形；幂，面积；讙，一说通“合”。}

\begin{questions}
\question[4] 解释材料一中画线的词：(1) 割之\underline{弥}细；(2) 与圆周\underline{合体}；(3) 又有\underline{余径}；(4) 因而\underline{三}之。
\answerspace{2cm}
\begin{solution}
(1) 弥：越，更加。(2) 合体：合为一体，重合。(3) 余径：正多边形一边之外、到圆周为止多出来的那段半径（半径减去边心距）。
(4) 三：数词用作动词，乘以三（三倍之）。每个 1 分。\score{4}
\end{solution}

\question[6] 把下面的句子翻译成现代汉语。
\begin{enumerate}
  \item 割之又割，以至于不可割，则与圆周合体而无所失矣。
  \item 以一面乘半径，觚而裁之，每辄自倍。
\end{enumerate}
\answerspace{3cm}
\begin{solution}
(1) 不断地分割下去，一直到不能再分割，（多边形）就与圆周合为一体，没有什么损失（误差）了。\score{3}
(2) 用（多边形的）一条边乘半径，按多边形的每条边把它裁分（成以圆心为顶点的小三角形），每个乘积都是相应小三角形面积的两倍。\score{6}
（“觚而裁之”译出按边裁分 1 分，“每辄自倍”译出“每一个都是两倍”1 分，语句通顺 1 分。）
\end{solution}

\question[4] 刘徽是怎样得出“以半周乘半径而为圆幂”（圆面积等于半周长乘半径）的？请用现代数学语言概括他的思路。
\answerspace{3cm}
\begin{solution}
用圆内接正多边形逼近圆：边数不断加倍，多边形面积与圆面积的差越来越小（“所失弥少”）；\score{1}
多边形每条边与圆周之间的“余径”给出误差的上界（“以面乘余径，则幂出弧表”），边数无限增加时余径趋于零，上下两界相夹；\score{3}
极限时多边形与圆重合，把多边形按边分成以圆心为顶点的小三角形（“觚而裁之”），面积为 $\frac12\times$ 周长 $\times$ 半径，即半周乘半径。
体现了极限与夹逼的思想。\score{4}
\end{solution}

\question[6] 材料二是《墨经》对几何概念的定义。
\begin{enumerate}
  \item 用现代数学语言解释“圜，一中同长也”，并说明它与今天圆的定义有何异同；
  \item 从“平，同高也”“中，同长也”“直，参也”中任选一条，说说它是用什么方式下定义的。
\end{enumerate}
\answerspace{3cm}
\begin{solution}
(1) 圆是有一个中心、其上各点到中心距离都相等的图形。\score{1}
与今天“平面内到定点的距离等于定长的点的集合”一致，抓住了“中心”与“等距”两个要素；\score{2}
但没有说明“平面”，没有用“点的集合”的说法，也没有区分圆周与圆面，表述较为简约。\score{3}
(2) 示例：“平，同高也”——用度量关系（各处高度相同）来规定“平”，是以性质、关系下定义；
“中，同长也”——中点到两端距离相等；“直，参也”——三点相参（对齐）即为直，用可检验的操作来定义。
说出所选条目的含义 2 分，说明定义方式（用关系、度量或可操作的检验来定义）1 分。\score{6}
\end{solution}
\end{questions}

% -----------------------------------------------------------------------------
\examsection*{语言文字运用}{本题共 1 小题，10 分}

\begin{questions}
\question[10] 下面这段话的四个句子各有一处表述不当，请指出并修改，使语意准确、简洁、通顺。

{\kaishu ①据统计，超过一半以上的同学认为数学证明很难。②其主要原因是因为他们缺乏系统的训练所造成的。
③事实上，通过反复练习，就能使证明能力得到提高。④证明不仅能检验结论是否正确，而且训练思维也离不开它。\par}
\answerspace{3cm}
\begin{solution}
①“超过”与“以上”语义重复，删去“以上”，改为“超过一半的同学”。\score{2.5}
②“原因是因为……所造成的”句式杂糅，改为“主要原因是他们缺乏系统的训练”。\score{5}
③“通过……使……”造成缺少主语，改为“只要反复练习，证明能力就能得到提高”。\score{7.5}
④前后陈述对象不一致，改为“证明不仅能检验结论是否正确，而且能训练思维”。\score{10}
\end{solution}
\end{questions}

% -----------------------------------------------------------------------------
\examsection*{写作}{本题 50 分}

\begin{questions}
\question[50] 阅读下面的材料，根据要求写作。

{\kaishu\hspace{2em}一位数学家说：“我宁可被一个好问题困住一个月，也不愿意马上被告诉答案。”

\hspace{2em}也有人说：“站在巨人的肩膀上才能看得更远。先把前人的结论接过来，才有余力去走前人没走过的路。”\par}

\hspace{2em}这两种说法引发了你怎样的思考？请结合你的学习经历与认识，写一篇文章。

\hspace{2em}要求：选准角度，确定立意，明确文体，自拟标题；不要套作，不得抄袭；不得泄露个人信息；不少于 800 字。
\answerspace{2cm}
\begin{solution}
\textbf{审题要点}　两种说法并不简单对立：前者强调亲身探索带来的理解与能力，后者强调继承与效率。好的立意应看到二者的张力与统一，
例如：接过结论之后仍要“重走一遍”使之成为自己的；或者区分不同阶段、不同目的下“困住”与“接过”的取舍。只赞一方、贬低另一方而不加分析的，
立意不超过二等。

\textbf{评分等级}（建议）
\begin{itemize}
  \item 一等（45—50 分）：立意深刻，能辨析两种说法的适用条件；论证有层次，例证具体（最好来自真实的学习经历）；语言准确流畅。
  \item 二等（36—44 分）：立意明确，有一定分析；结构完整；语言通顺。
  \item 三等（27—35 分）：立意基本符合材料，分析较浅，例证一般化。
  \item 四等（26 分以下）：偏离材料，或明显套作，或字数严重不足。
\end{itemize}
\score{50}
\end{solution}
\end{questions}

\end{paper}
\end{document}
